\section{Implementation and Costs}\label{sec:costs}\label{sec:impl}\label{sec:params}

\leanseal implements \cref{sec:construction} as a versioned format with a written specification, a Rust implementation and command-line tool, and an independent Python implementation that agrees with it on every test vector~\cite{leanSealRepo}. The solver saves its progress, so an interrupted solve can be resumed, and it can solve several puzzles at once, one per core. The Rust code depends only on \texttt{getrandom} and \texttt{signal-hook}, uses the ARMv8 and x86 SHA-256 instructions when present, and is tested with known answers, round trips, property tests over 480,000 mutated encodings and fuzzing of every decoder. The specification and an earlier C implementation were reviewed over several rounds by an AI model; no human cryptographer has reviewed the code or the format.

\paragraph{Encryption.}
Encryption uses only SHA-256. For a fresh 16-byte nonce, subkeys $k_{\mathrm{enc}}$ and $k_{\mathrm{mac}}$ are labelled hashes of $K$ and the nonce; the keystream is $\HMAC(k_{\mathrm{enc}}, \mathsf{be64}(j))$ for block $j$, and an HMAC tag under $k_{\mathrm{mac}}$ covers every preceding byte. A sealed transaction of $L$ bytes takes $128 + nb + L$ bytes: the $68 + nb$-byte puzzle, 8 bytes of framing, the nonce and length (20 bytes) and the tag (32 bytes). One puzzle can seal several messages under different nonces; they all open together.

\subsection{Measurements}\label{sec:measurements}

\Cref{tab:measurements} gives the measurements on two Apple laptops: a fanless MacBook Air with an M4 (10 CPU cores) and a MacBook Pro with an M4 Max (16 CPU cores). GPU figures come from a separate Metal benchmark of the chain step; the tool does not yet build puzzles on a GPU.

\begin{table}[t]
  \centering\small
  \caption{Measurements. The measured puzzle is a one-minute lock with $n = 16$, $m = 2^{30}$ and 24-byte links; the 30-second reference puzzle ($m = 2^{29}$) takes half these times to build and to open without the key. The M4 Air building and opening runs used 32-byte links, with the same chain work. Opening with the key is the command-line tool's end-to-end time, including process start and file input and output.}
  \label{tab:measurements}
  \begin{tabularx}{\textwidth}{@{}>{\raggedright\arraybackslash}Xrr@{}}
    \toprule
    \colhead{Measured} & \colhead{M4 Air} & \colhead{M4 Max} \\
    \midrule
    One chain step, one CPU core & 23.2\,ns & 24.1\,ns \\
    Total CPU hash rate, independent chains & 461\,M/s & 867\,M/s \\
    Total GPU hash rate, many chains in flight & 0.71\,G/s & 2.96\,G/s \\
    Building the one-minute puzzle on the CPU & 42.1\,s & 30.1\,s \\
    Opening it without the key, one core & 6\,min 42\,s & -- \\
    Opening a sealed 110-byte transaction with the key & 12.0\,ms & -- \\
    \bottomrule
  \end{tabularx}
\end{table}

\subsection{What a delay costs}\label{sec:building}

Four formulas describe the scheme. The delay is $T = N\ell$. Building takes about $N / \Rs$, where $\Rs$ is the sender's total hash rate. The puzzle takes $68 + nb$ bytes. Opening without the key takes $N$ times the opener's own time per step: 3.3 minutes on one laptop core for a 30-second puzzle.

The sender's cost is therefore proportional to the delay: about 0.70 seconds of a laptop's time, with all its cores busy, per second of delay (0.097 on the M4 Max GPU, amortized over a batch that takes hours). For constructions from a random oracle, this work cannot be reduced asymptotically~\cite{MMV11} (\cref{sec:limits}). Two observations make this cost acceptable. First, a puzzle does not depend on the transaction it will seal, so a wallet can keep a stock and build more while idle or charging; the building time is then a cost, not a delay. A laptop CPU running all day builds about 2,800 30-second puzzles (\cref{sec:chains}). Second, the delay only needs to cover the time until the ordering is fixed (\cref{sec:timing}); halving the delay halves all the work and leaves the puzzle size unchanged.

Precomputation removes the latency, not the cost. Building is spent on every sealed transaction, including those revealed on time, unlike the solvers' work, which arises only on a missed reveal (\cref{sec:solvers}). At the 265 transactions per slot of our mainnet sample (\cref{sec:size}), one fresh 30-second puzzle each keeps about 460 M4 Air laptops fully busy at 21\,s per puzzle, or about 680 at the stock rate below. At about 20\,W per laptop under load, an estimate we have not measured, this is 0.1 to 0.2\,Wh per transaction and 9 to 14\,kW in total. We have not measured phones; a wallet on a phone would build its stock while charging, or outsource building (\cref{sec:sharing}).

\subsection{The number of chains}\label{sec:chains}

The delay depends only on $N$; the number of chains $n$ sets how many cores can work on one puzzle at once, and the size. A sender that builds a puzzle when its transaction is ready needs one chain per core it wants to use: $n = 16$ builds a 30-second puzzle in about 21\,s on the M4 Air and 15\,s on the M4 Max, half the measured times for a one-minute puzzle. A sender that builds a stock obtains its parallelism from building many puzzles at once instead (five one-minute 16-chain puzzles built together on the M4 Air finished in 187.6\,s, or 37.5\,s amortized per puzzle, about 19\,s per 30-second puzzle), so it can use a single chain per puzzle. \Cref{tab:chains} compares the options.

\begin{table}[t]
  \centering\small
  \caption{30-second puzzles ($N = 2^{33}$, 24-byte links) for different numbers of chains. The times are half those measured for one-minute puzzles, except for $n = 4$, which is projected from the measured time per step.}
  \label{tab:chains}
  \begin{tabular}{@{}rrrl@{}}
    \toprule
    \colhead{Chains $n$} & \colhead{Puzzle} & \colhead{Added to a transaction} & \colhead{Building on a 10-core laptop} \\
    \midrule
    16 & 452\,B & 512\,B & on demand, about 21\,s \\
    4 & 164\,B & 224\,B & on demand, about 50\,s \\
    1 & 92\,B & 152\,B & from a stock: one per 30\,s, about 2,800 a day \\
    \bottomrule
  \end{tabular}
\end{table}

With $n = 1$, sealing a 110-byte transfer gives 262 bytes rather than 622, and the guessing term of \cref{thm:delay} vanishes. The daily figure above is projected from a single run, which was not archived and was made on a machine that was not idle: ten single-chain puzzles built at once on the M4 Air, one per core, took 75\,s for $2^{31}$ steps each. A stock builder could interleave two single-chain puzzles per core, as the builder already does with chains, to keep the 1.5-fold gain of two SHA-256 pipelines per core.
